% =====================================================================
%  Introduction générale
% =====================================================================
\chapter*{Introduction générale}
\addcontentsline{toc}{chapter}{Introduction générale}
\markboth{Introduction générale}{Introduction générale}

La communication constitue aujourd'hui un pilier stratégique de toute
organisation. Qu'il s'agisse d'une campagne promotionnelle par e-mail,
d'une facture adressée à un client, d'un contrat à signer ou d'un message
enrichi sur mobile, chaque point de contact véhicule l'image de marque de
l'entreprise et participe à la relation qu'elle entretient avec ses
publics. La multiplication des canaux numériques a rendu cette
communication à la fois plus riche et plus exigeante : les entreprises
doivent produire, rapidement et à grande échelle, des supports cohérents,
personnalisés et fidèles à leur identité visuelle.

Or la fabrication de ces supports demeure, dans bien des cas, un exercice
lent et technique. Concevoir un e-mail responsive correct suppose de
maîtriser des langages de balisage tels que le HTML ou le MJML~; produire
une facture ou un contrat exige de la rigueur documentaire~; et l'ensemble
requiert le plus souvent l'intervention conjointe de profils créatifs et
techniques. Les équipes métier, en particulier les équipes marketing,
souhaitent au contraire pouvoir \emph{décrire} ce qu'elles veulent — ou
\emph{réutiliser} un visuel existant — et obtenir instantanément un
résultat professionnel et immédiatement modifiable.

C'est dans ce contexte que s'inscrit le présent projet de fin d'études,
réalisé au sein de \textbf{France Téléphone}, filiale du groupe
\textbf{Winvest Capital}. Il a consisté à concevoir et à développer
\textbf{Winaity Template Builder}, une plateforme logicielle en mode SaaS
(\emph{Software as a Service}) multi-entreprises destinée à la création, à
la gestion et à la génération de modèles de communication sur cinq canaux :
l'e-mail, la facture, le contrat, le SMS et le message RCS. Le c\oe{}ur de
la plateforme est un éditeur visuel par glisser-déposer, augmenté d'un
assistant d'intelligence artificielle \emph{agentique} capable de générer
et de modifier de véritables blocs éditables, voire de reconstruire une
campagne complète à partir d'une simple affiche marketing.

L'avènement récent des grands modèles de langage a ouvert des perspectives
inédites en matière d'assistance à la création de contenu. Le projet
explore précisément cette convergence entre l'ingénierie logicielle et
l'intelligence artificielle, en s'appuyant sur un modèle de langage
auto-hébergé afin de préserver la souveraineté des données de l'entreprise.

Ce rapport rend compte de la démarche suivie et des réalisations
accomplies. Il s'organise en six chapitres.

Le \textbf{premier chapitre} présente le cadre général du projet :
l'organisme d'accueil, la problématique, une étude critique des solutions
existantes, la solution proposée et la méthodologie de travail retenue.

Le \textbf{deuxième chapitre} est consacré à l'analyse des besoins et à la
conception : identification des acteurs, expression des besoins
fonctionnels et non fonctionnels, diagrammes de conception, architecture de
l'application, backlog produit et environnement de développement.

Les \textbf{quatre chapitres suivants} détaillent le déroulement des quatre
sprints du projet. Le troisième chapitre traite de l'authentification, de
l'isolation multi-entreprises et de la gestion des utilisateurs. Le
quatrième porte sur l'éditeur de modèles multi-canal. Le cinquième, c\oe{}ur
de la valeur ajoutée du projet, décrit l'assistant d'intelligence
artificielle agentique. Le sixième aborde les abonnements, les
intégrations et la mise en production.

Enfin, une \textbf{conclusion générale} dresse le bilan du travail réalisé,
des compétences acquises et des perspectives d'évolution de la plateforme.
